فرض کنید $X\in\mathbb{R}^{n\times d}$ و $X^T X = nI$ و ستون $i$ام برابر $x_{*i}\in\mathbb{R}^{n}$ است.

\subsection*{۱}
\begin{flalign*}
J_1(w)
&= \|Xw-y\|^2 + \lambda\|w\|_1 \\
&= (Xw-y)^T(Xw-y)+\lambda\sum_{i=1}^d |w_i| \\
&= w^T X^T X w -2y^T Xw + \|y\|^2 + \lambda\sum_{i=1}^d |w_i| \\
&= n\|w\|^2 -2\sum_{i=1}^d w_i\, x_{*i}^T y + \|y\|^2 + \lambda\sum_{i=1}^d |w_i| \\
&= \|y\|^2 + \sum_{i=1}^d \Big( n w_i^2 -2 w_i (x_{*i}^T y) + \lambda |w_i| \Big). &&
\end{flalign*}
پس می‌توان نوشت:
\begin{flalign*}
J_1(w)=\|y\|^2+\sum_{i=1}^d f(x_{*i},w_i),
\qquad
f(x_{*i},w_i)\triangleq n w_i^2 -2 w_i (x_{*i}^T y) + \lambda |w_i|. &&
\end{flalign*}

\subsection*{۲}
به‌علت جدایی‌پذیری، برای هر $i$ کافی است
\begin{flalign*}
\phi_i(w_i)=n w_i^2 -2c_i w_i + \lambda|w_i|,
\qquad c_i\triangleq x_{*i}^T y
&&
\end{flalign*}
را کمینه کنیم. شرط لازم و کافی بهینگی (شرط زیرگرادیان):
\begin{flalign*}
0 \in \partial \phi_i(w_i)
= 2n w_i -2c_i + \lambda\,\partial |w_i|. &&
\end{flalign*}
سه حالت:

\begin{itemize}
\item اگر $w_i^*>0$، آنگاه $\partial|w_i^*|=\{1\}$ و:
\begin{flalign*}
2n w_i^* -2c_i + \lambda = 0
\;\;\Rightarrow\;\;
w_i^*=\frac{c_i-\lambda/2}{n},
\qquad
w_i^*>0 \Longleftrightarrow c_i>\lambda/2. &&
\end{flalign*}

\item اگر $w_i^*<0$، آنگاه $\partial|w_i^*|=\{-1\}$ و:
\begin{flalign*}
2n w_i^* -2c_i - \lambda = 0
\;\;\Rightarrow\;\;
w_i^*=\frac{c_i+\lambda/2}{n},
\qquad
w_i^*<0 \Longleftrightarrow c_i<-\lambda/2. &&
\end{flalign*}

\item اگر $w_i^*=0$، آنگاه $\partial|0|=[-1,1]$ و:
\begin{flalign*}
0 \in -2c_i + \lambda[-1,1]
\;\;\Longleftrightarrow\;\;
|2c_i|\le \lambda
\;\;\Longleftrightarrow\;\;
|c_i|\le \lambda/2. &&
\end{flalign*}
\end{itemize}

\subsection*{۳}
\begin{flalign*}
J_2(w)
&=\|Xw-y\|^2+\lambda\|w\|_2^2
= w^T X^T X w -2y^T Xw + \|y\|^2 + \lambda w^T w \\
&= (n+\lambda)\|w\|^2 -2\sum_{i=1}^d c_i w_i + \|y\|^2. &&
\end{flalign*}
شرط بهینگی:
\begin{flalign*}
\nabla J_2(w)=2(n+\lambda)w -2X^T y = 0
\;\;\Rightarrow\;\;
w^\#=\frac{1}{n+\lambda}X^T y. &&
\end{flalign*}
پس:
\begin{flalign*}
w_i^\#=0
\;\;\Longleftrightarrow\;\;
(x_{*i}^T y)=0. &&
\end{flalign*}

\subsection*{۴}
طبق بخش (۲)، در \lr{Lasso} اگر $|x_{*i}^T y|\le \lambda/2$ باشد، دقیقاً $w_i^*=0$ می‌شود (آستانه‌گذاری و ایجاد صفر).
اما طبق بخش (۳)، در \lr{Ridge} داریم $w_i^\#= \frac{x_{*i}^T y}{n+\lambda}$ و فقط وقتی صفر می‌شود که دقیقاً $x_{*i}^T y=0$ باشد.
بنابراین معمولاً $w^*$ تُنُک‌تر از $w^\#$ است.